\documentclass[10pt,a4paper]{amsart}
\usepackage[margin=19mm]{geometry}
\usepackage{amsmath,amssymb,mathtools}
\usepackage[scr=rsfs,cal=euler]{mathalfa}
\usepackage{xfrac}
\usepackage[unicode,hidelinks,bookmarks=false]{hyperref}
\newcommand{\ie}{i.e.\ }
\newcommand{\cf}{cf.\ }
\newcommand{\resp}{respectively\ }
\newcommand{\Subsection}{\subsection}
\providecommand{\nobreakdash}{\nobreak\hbox{-}}
\setlength{\emergencystretch}{2em}
\multlinegap0pt
\usepackage{amssymb}
\usepackage{mathtools}
\usepackage{csquotes}
\usepackage{algorithm}\usepackage{algpseudocode}
\usepackage{tikz}
\usepackage{enumerate}
\newtheorem{lem}{Lemma}
\newtheorem{thm}{Theorem}
\newcommand{\fac}{\hbox{\rm Fac}}
\newcommand{\pri}{\mathrm{Prim}}
\newcommand{\rp}{\mathrm{RP}}
\title{A Tighter Upper Bound for the Number of Distinct Squares in Circular Words}
\author{Shuo Li, Yuan Song}
\date{Source: arXiv:2605.12215v1 (CC BY 4.0)}
\begin{document}
\maketitle

\noindent\emph{Source and licence.} A Tighter Upper Bound for the Number of Distinct Squares in Circular Words. Version arXiv:2605.12215v1 (CC BY 4.0), distributed by arXiv under the Creative Commons Attribution 4.0 International licence, http://creativecommons.org/licenses/by/4.0/, as recorded in the arXiv licence metadata for this exact version and shown on the abstract page https://arxiv.org/abs/2605.12215v1. Full licence notice: \texttt{licenses/arxiv-2605.12215-CC-BY-4.0.txt}.\par\smallskip

\noindent\textit{Editorial scope.} Counted unit: the article's thm thm:rp-upper-bound (source0.tex lines 2035--2041). The article's complete original proof of it is delivered with it (source0.tex lines 2043--2372, one \texttt{proof} environment of 330 lines). No mathematics is written, repaired or re-derived: the statement and the proof are the article's own lines, in the article's own notation, and no retained line is substituted or re-flowed.  Retained uncounted as the source-local context the proof needs: lines 1098--1103; lines 1149--1158; lines 1219--1232; lines 1279--1288; lines 1383--1391; lines 1528--1553; lines 1755--1781.  Nothing was omitted from any retained span, so the package carries no omission record.  No retained line contains a citation, so the wrapper carries no bibliography and no bibliography entry.  No retained line contains a figure environment, an image-inclusion command or drawing code, so no image is delivered and the package declares no asset dependency.  Editorial formatting only, in addition: an AMS \texttt{amsart} preamble that restates the article's own declarations and macros used by the retained lines, namely the environments \texttt{lem}, \texttt{thm} on separate counters, together with the standard packages those lines need.  The retained environments print in this document as Lemma 1, Theorem 1 in order of appearance, whereas the article prints the same items with the numbering of its own full document.  The complete native source of the article as distributed in the arXiv e-print for this exact version is bundled unchanged as \texttt{sources/200409/source0.tex}.
\begin{lem}\label{restriction on length and number of sc}
For every $u\in\pri(w)$ such that $u^2\in\fac(w)$, one has
\[
2|u|+M(u)-1\le n.
\]
\end{lem}

\begin{lem}\label{lem:A-optimization}
Let
\[
        c_0:=2-\sqrt3.
\]
If \(\ell+m\le A\), then
\[
        F(\ell,m)\le c_0Am.
\]
\end{lem}

\begin{lem}\label{lem:B-optimization}
Let
\[
        c_*:=\frac{3-\sqrt5}{4}.
\]
If
\[
        2\ell+m\le B,
\]
then
\[
        F(\ell,m)\le c_*Bm.
\]
\end{lem}

\begin{lem}\label{lem:pairwise-A}
Let \(i<j\). If
\[
        m_j\ge \ell_i,
\]
then
\[
        2\ell_j+m_j+\ell_i\le n+1.
\]
\end{lem}

\begin{lem}\label{lem:large-large-unique}
For all sufficiently large \(n\), there is at most one conjugacy class
\([u]\in D(w)\) such that
\[
        |u|>\frac n4,
        \qquad
        M(u)>\frac n4 .
\]
\end{lem}

\begin{lem}\label{lem:optimization-case1}
Let
\[
c_0:=2-\sqrt{3},\qquad c^*:=\frac{3-\sqrt{5}}{4}.
\]
For \(0\leq a\leq \frac14\) and \(0\leq b\leq 1-2a\), define
\[
E(a,b):=
\max\left\{
\frac a2,\,
\min\{(c_0(a+b),c^*(1-a)\}
\right\}.
\]
Define
\[
\Phi_1(a,b):=
\begin{cases}
\dfrac12 b^2+E(a,b)(1-b), & b<a,\\[6pt]
ab-\dfrac12a^2+E(a,b)(1-b), & b\geq a.
\end{cases}
\]
Then
\[
\Phi_1(a,b)\leq \frac{95-25\sqrt5}{236}.
\]
\end{lem}

\begin{lem}\label{lem:optimization-case3}
Let
\[
c^*:=\frac{3-\sqrt5}{4}.
\]
For \(\frac14<a<\frac38\), define
\[
K(a):=
\max\left\{
\frac18,\,
\frac14-\frac1{32a},\,
(1-a)c^*
\right\}.
\]
For \(b\in\left(\frac14,1-2a\right]\), define
\[
\Phi_3(a,b):=
\begin{cases}
\dfrac12 b^2+K(a)(1-b), & b<a,\\[6pt]
ab-\dfrac12a^2+K(a)(1-b), & b\geq a.
\end{cases}
\]
Then
\[
\Phi_3(a,b)\leq \frac{3(4-\sqrt5)}{32}.
\]
\end{lem}

\begin{thm}\label{thm:rp-upper-bound}
Let \(w\) be a word of length \(n\). Then
\[
        \rp(w)\le
        \left(\frac{95-25\sqrt5}{236}\right)n^2+O(n).
\]
\end{thm}

\begin{proof}
Set
\[
        a:=\frac{\ell_1}{n},
        \qquad
        b:=\frac{m_1}{n}.
\]
By Lemma~\ref{restriction on length and number of sc},
\[
        2a+b\le 1+\frac1n.
\]
In the following optimization estimates, replacing this constraint by
\(2a+b\le1\) only changes the final estimate by \(O(n)\). Indeed, all the
functions involved are piecewise quadratic on compact regions, hence are
Lipschitz on a fixed neighbourhood of these regions.

We divide the proof into three cases.

\medskip

\noindent
\textbf{Case 1: \(\ell_1\le \frac n4\).}

Then
\[
        0\le a\le\frac14,
        \qquad
        0\le b\le 1-2a
\]
up to an \(O(1/n)\)-error, which is absorbed into the final \(O(n)\)-term.

The first class contributes at most
\[
        F(\ell_1,m_1)
        \le
        n^2
        \begin{cases}
        \dfrac12b^2,& b<a,\\[1mm]
        ab-\dfrac12a^2,& b\ge a.
        \end{cases}
\]

Now let \(i\ge2\). We estimate \(F(\ell_i,m_i)\).

If
\[
        m_i<\ell_1,
\]
then
\[
        F(\ell_i,m_i)\le \frac12m_i^2
        \le
        \frac12\ell_1m_i
        =
        \frac a2\,n\,m_i.
\]

If
\[
        m_i\ge\ell_1,
\]
then we have two independent estimates. First, since \(A_i\le A_1\),
Lemma~\ref{lem:A-optimization} gives
\[
        F(\ell_i,m_i)
        \le
        c_0A_1m_i
        =
        c_0(a+b)n\,m_i.
\]
Second, Lemma~\ref{lem:pairwise-A} gives
\[
        2\ell_i+m_i+\ell_1\le n+1.
\]
Hence
\[
        2\ell_i+m_i\le n+1-\ell_1.
\]
By Lemma~\ref{lem:B-optimization},
\[
        F(\ell_i,m_i)
        \le
        c_*(n+1-\ell_1)m_i
        =
        c_*(1-a)n\,m_i+O(m_i).
\]
Therefore, when \(m_i\ge\ell_1\),
\[
        F(\ell_i,m_i)
        \le
        \min\{c_0(a+b),c_*(1-a)\}\,n\,m_i+O(m_i).
\]

Combining the two cases, for every \(i\ge2\),
\[
        F(\ell_i,m_i)
        \le
        E(a,b)n\,m_i+O(m_i),
\]
Since
\[
        \sum_{i=2}^t m_i\le n-m_1=(1-b)n,
\]
we get
\[
        \sum_{i=1}^t F(\ell_i,m_i)
        \le
        n^2\Phi_1(a,b)+O(n).
\]
By Lemma~\ref{lem:optimization-case1},
\[
        \sum_{i=1}^t F(\ell_i,m_i)
        \le
        \left(\frac{95-25\sqrt5}{236}\right)n^2+O(n).
\]

\medskip

\noindent
\textbf{Case 2: \(\ell_1>\frac n4\) and \(m_1\le\frac n4\).}

Then
\[
        a>\frac14,
        \qquad
        b\le\frac14.
\]
Since \(b\le a\), we have
\[
        F(\ell_1,m_1)=\frac12m_1^2=\frac12b^2n^2.
\]

For \(i\ge2\), since \(A_i\le A_1\), Lemma~\ref{lem:A-optimization} gives
\[
        F(\ell_i,m_i)\le c_0A_1m_i
        =
        c_0(a+b)n\,m_i.
\]
Hence
\[
        \sum_{i=1}^tF(\ell_i,m_i)
        \le
        n^2\left[
        \frac12b^2+c_0(a+b)(1-b)
        \right]+O(n).
\]
For fixed \(b\), the expression in brackets is increasing in \(a\). By
\(2a+b\le1\), its maximum occurs at
\[
        a=\frac{1-b}{2}.
\]
Therefore
\[
\begin{aligned}
        \frac12b^2+c_0(a+b)(1-b)
        &\le
        \frac12b^2+\frac{c_0}{2}(1-b^2)\\
        &=
        \frac{c_0}{2}+\frac{1-c_0}{2}b^2.
\end{aligned}
\]
Since \(0\le b\le1/4\), the last expression is at most
\[
        \frac{c_0}{2}+\frac{1-c_0}{32}
        =
        \frac{31-15\sqrt3}{32}.
\]
Thus
\[
        \sum_{i=1}^tF(\ell_i,m_i)
        \le
        \frac{31-15\sqrt3}{32}n^2+O(n).
\]
Finally,
\[
        \frac{31-15\sqrt3}{32}
        <
        \frac{95-25\sqrt5}{236}.
\]
Hence Case 2 is bounded by the desired constant.

\medskip

\noindent
\textbf{Case 3: \(\ell_1>\frac n4\) and \(m_1>\frac n4\).}

By Lemma~\ref{lem:large-large-unique}, \([u_1]\) is the only class satisfying
both
\[
        \ell_i>\frac n4,
        \qquad
        m_i>\frac n4.
\]
Thus for every \(i\ge2\), either
\[
        \ell_i\le\frac n4
        \quad\text{or}\quad
        m_i\le\frac n4.
\]

Since
\[
        a>\frac14,
        \qquad
        b>\frac14,
        \qquad
        2a+b\le1,
\]
we have
\[
        \frac14<a<\frac38.
\]

We claim that for every \(i\ge2\),
\[
        F(\ell_i,m_i)\le K(a)n\,m_i+O(m_i),
\]
where
\[
K(a):=
\max\left\{
        \frac18,\,
        \frac14-\frac{1}{32a},\,
        (1-a)c_*
\right\}.
\]

Indeed, if
\[
        m_i\le\frac n4,
\]
then
\[
        F(\ell_i,m_i)\le \frac12m_i^2\le \frac18n\,m_i.
\]

Now assume
\[
        m_i>\frac n4.
\]
Then necessarily
\[
        \ell_i\le\frac n4.
\]
If, in addition,
\[
        m_i<\ell_1=an,
\]
then
\[
        F(\ell_i,m_i)
        \le
        \frac n4m_i-\frac12\left(\frac n4\right)^2
        =
        \frac14n\,m_i-\frac1{32}n^2.
\]
Since \(m_i<an\), we have
\[
        -\frac1{32}n^2
        \le
        -\frac{1}{32a}n\,m_i+O(m_i).
\]
Therefore
\[
        F(\ell_i,m_i)
        \le
        \left(\frac14-\frac{1}{32a}\right)n\,m_i+O(m_i).
\]

Finally, if
\[
        m_i\ge \ell_1,
\]
then Lemma~\ref{lem:pairwise-A} gives
\[
        2\ell_i+m_i+\ell_1\le n+1.
\]
Hence
\[
        2\ell_i+m_i\le n+1-\ell_1.
\]
By Lemma~\ref{lem:B-optimization},
\[
        F(\ell_i,m_i)
        \le
        c_*(n+1-\ell_1)m_i
        =
        (1-a)c_*\,n\,m_i+O(m_i).
\]
This proves the claim.

Since
\[
        \sum_{i=2}^tm_i\le n-m_1=(1-b)n,
\]
we get
\[
        \sum_{i=1}^tF(\ell_i,m_i)
        \le
        n^2\Phi_3(a,b)+O(n).
\]
By Lemma~\ref{lem:optimization-case3},
\[
        \sum_{i=1}^tF(\ell_i,m_i)
        \le
        \frac{3(4-\sqrt5)}{32}n^2+O(n).
\]
Since
\[
        \frac{3(4-\sqrt5)}{32}
        <
        \frac{95-25\sqrt5}{236},
\]
Case 3 is also bounded by the desired constant.

\medskip

The three cases cover all possibilities. Therefore
\[
        \sum_{i=1}^tF(\ell_i,m_i)
        \le
        \left(\frac{95-25\sqrt5}{236}\right)n^2+O(n).
\]
Consequently,
\[
        \rp(w)
        \le
        \left(\frac{95-25\sqrt5}{236}\right)n^2+O(n).
\]
\end{proof}

\end{document}
